\documentclass[man,12pt,natbib]{apa6}
\usepackage[colorlinks=false]{hyperref}
\usepackage{amssymb,amsmath,times,ragged2e,soul,textcomp,graphicx}
\linespread{1.5}

\begin{document}

\title{Homework 1: Computers as a New Species of Language-Capable Beings}
\shorttitle{Homework 1}
\author{Edward Hern\'{a}ndez}
\date{12 February}
\affiliation{College of William \& Mary}
\maketitle

\emph{Will computers ever achieve the ability to produce and understand human
language?}
To answer this seemingly simple question, I will break it into three of its
(many) component questions. Can a computer produce language? Can a computer
understand language?  If it can, will the language that it can produce or
understand be \emph{human}?

For the purposes of this paper, I will assume that any breakthrough which is
conceptually possible will be eventually technologically realized. If a
computer can, in principle, produce and understand language, I will assume that
someone\footnote{\st{Google.} Alphabet.} will eventually build a computer which
does.  I also assume that the computational theory of mind (termed ``strong
AI'' by \citealp{Searle80}) falls to (at least) embodied cognition
arguments.\footnote{ I leave open the possibility that CTM is not false
(specifically that it is trivially true and compatible with embodied
cognition). However, if I held it to be true, the answer to the question would
be trivial, and the focus of any paper on the topic would be to prove CTM. In
order to write a topical essay, I caveat embodied cognition is true and
incompatible with CTM.} If it is the case that the human brain (or its mind) is
a computer, the above question is moot, as human language has then never been
produced or understood by anything but computers, and building non-human
computers which understand and produce language by the very same computational
processes is fully possible and will inevitably be accomplished.

First, can a computer produce language?
Yes. It can.
This does not need to be justified conceptually, because computers already do
produce language.  Computers routinely produce outputs in formal languages
(like assembly languages). They also produce strings of language which are
understandable by humans, like error messages or application feedback. These
outputs can be printed to screens or pages, or spoken aloud.  They also produce
novel strings of language,\footnote{I shy away from saying ``in a language''
because of the superdiversity literature I discuss in more detail below.} like
haiku \citep{Cope11}, or conversations \citep{Biever12,Baraniuk14}.

Can a computer understand language?
This is a much larger question, and requires some clarification before it can
be answered. What does it mean to understand? Does producing language
necessitate understanding it?

\citet{Searle80} explores exactly these questions. To use his thought
experiment, imagine a person who does not speak Chinese in a room with books of
Chinese. That person is passed notes, written in Chinese, through a door, and
writes notes in response based on the characters in the books. She has no idea
what the characters on the notes or her responses mean (to the people outside
the room), but eventually she learns to use the books so well that everyone
outside the room believes her to be a fluent (human) speaker of Chinese. In
this case we are inclined to say that the person does not \emph{understand}
Chinese or the language which constitutes her inputs and outputs.

This corresponds directly to the question of whether computers understand
language. If we say of the person in the Chinese Room that she does not
understand the notes, then we must say of a computer that it does not
understand language which it produces based on machine-learning, even if it
passes other measures, like the Turing test. The human, after all, would also
pass as a fluent speaker of Chinese, and this does not move us intuitively to
say she understands.

\citet{Searle80} takes this thought experiment to show us a divide based on
intentionality. Brains, he argued, posses special causal capacity, enabling
them to be intentional in a way which a computer (or, other computers, on his
view) cannot.  Another possible explanation of our reticence is that the
operations within the room are (arguably) exclusively syntactic
\citep{Marconi96}. Words are chosen from books, fit together such that they
adhere to the rules of Chinese and returned. The semantic content of both
incoming and outgoing notes is completely opaque to the person writing the
responses.  Both of these are robust attacks, I here argue that they are
incorrect about it is to understand.

Imagine that you are asked whether a given human understands language. How do
you come to an answer?  If you are like me, you ask them questions, and when
they answer in correct and comprehensible ways, you conclude that they do.
This means that computers which pass the Turing test \emph{do} understand
language by our everyday definition of the word understand. I hold that this is
entirely sufficient to say that they do understand, based on arguments out of
\citet{Wittgenstein53,Wittgenstein69}.  There is no language game for
determining understanding which involves assessing the causal effect of that
language on the brain or whether that brain engages with the semantic content
rather than solely the words' syntactic information. Obviously this cannot be
our criteria for other humans, because that data is never available to us as
interlocutors. If a human can talk about something, I judge that she
understands both language and the semantic content contained in our utterances.
Thus neither their supposed lack of intentionality or alleged failure to engage
with semantic content, computers should be taken to understanding language if
we would say of them that they understood language were they human. Computers
understand language (and anything else) in that they satisfy the criteria in
our language-games for understanding.

Additionally, since I am operating under the assumption that all possible
technological advances \emph{will} be made, I expect that biological computers
(which are possible in at least principle) will be produced such that they have
any possible biological property. I see no reason that such a computer would be
theoretically incapable of intentionality in the way that Searle requires, and
might store and access semantic information in ways (functionally, if not
physically) identical to brains, satisfying the other concern.

Lastly, is the language which computers understand or produce human? This may
seem like a trivial question to ask, but it does not fall within the scope of
either of the above component questions, and is a requirement of the overall
question. I here argue that, currently, computers do not produce or understand
human language, and that there are significant conceptual and technical hurdles
to their doing so.

Currently, computers produce a small subset of the language that is created by
humans, and they are very rarely judged to be human speakers. Recently, the
first computer arguably passed the Turing test by convincing 33\% of judges
that it was human \citep{Reading14}. By my above argumentation, that computer
produced and understood language. However, it was only able to do so in a
specific context: a conversation via keyboard in which the other interlocutor
believed it to be a 13 year old Ukranian boy. This ruse was designed to cover
up the computer's errors and failures. Those failures were sufficient that we
might well say of a human who made them did not understand language (well).
Computers also routinely fail to understand language on speech-to-text and text
tasks, fail to respond correctly to sarcasm and metaphor, and are horribly bad
at handling dialectal variation. Until these hurdles are overcome, none of
these capabilities constitute production or understanding human language.

There are, of course, more advanced computers targeted at approximating human
language. Generally, linguists or other language scientists attempt to run
machine-learning programs on these computers.  However they feed computers
``language'' input which does not closely resemble language as it is used
socially. It is often downcased, sometimes with punctuation removed. Spelling
is standardized. Only forms deemed grammatical in the target language are
selected for the input. This is all based on the practices of Chomskyan
linguists, who hold that the ideal speaker in a homogeneous speech community is
a reasonable model for the language learner (and therefore a language-capable
computer). Because this is their model for language, they feed computers
homogeneous strings judged to be grammatical in a single language.

However, no speech community is homogeneous. Women speak differently than men,
and both speak differently at different ages. There are generational and class
differences among speakers. There are also, of course, multilingual speakers,
whose communities are obviously not homogeneous. In many places, there are
multilingual communities, in which nearly all communication may be
heterogeneous and heteroglossic, as the text in Figure 1, which is written
ungrammatically and simultaneously in multiple dialects of Chinese, with a mix
of two incompatible styles of Chinese characters.  These realities have put
under threat the very idea of ``a language,'' moving some, like
\citet{Blommaert11} to produce and adopt new models of language as a whole. On
these views, often broadly labeled \emph{superdiversity}, language is much less
predictable than Chomskyan linguists assume \citep[p.~1]{Blommaert11}, varying
on every level, based on not only the interlocutors and their languages, but
also based on every sort of context available.

\begin{figure}
	\resizebox{\columnwidth}{!}{
	\includegraphics{./.resources/antwerp.png}
}
	\caption{A note in a shop window in Antwerp. \textcopyright~Jan Blommaert}
\end{figure}

Additionally, neither the speech signal nor visual input is clean and in
formats that computers can understand. The speech signal is never alone. There
are ambient sounds that must be filtered out or ignored. There is phonetic
variation which is not semantically relevant. The visual field is perhaps more
fraught. Language comes in any number of typefaces and even more varieties of
handwriting. Figure 1 shows language in multiple dialects and multiple forms of
the same hanzi. Figure 2 shows an advertisement on a crowded bulletin board.
Most of the language on the board is recycled from older sources, and are,
apart from being somewhat difficult to read, semantically dependent on their
having been transplanted. How could a computer take this image and know from
the appearance of the language that it originates elsewhere (which I firmly
believe is part of understanding the language present)?

\begin{figure}
	\resizebox{\columnwidth}{!}{
		\includegraphics{./.resources/bulletin.jpg}
	}
	\caption{An advertisement on a bulletin board}
\end{figure}

Computers are very far from being able to understand (much less respond
appropriately to and themselves produce) such variation.  All of the above
challenges can eventually be overcome with bigger and better training sets for
conventional computers, but, according to the embodied cognition literature,
much of human language processing is handled by the interaction of the
sensorimotor capacities of the brain with the body and the environment. This
view also holds that understanding semantic content may involve sensorimotor
processing (e.g.  the concept of kicking may involve the sensorimotor activity
related to performing kicking motions; \citealp{Rueschemeyer10}).  This
conception of understanding is not contradictory with my Wittgensteinian
definition above. To assess whether someone understands some semantic content,
I might ask them to demonstrate it.  I might ask someone to kick, etc., and I
might judge that a person who fails to do so does not understand kicking. If a
computer has no sensorimotor functions related to kicking, then, it may be
judged to not understand it.\footnote{But, a robot which could both talk and
kick, and could talk about kicking, I would judge to understand both kicking
and language.}

Quite a lot of language as it is spoken by humans has to do with human
activities in which a computer cannot (yet) participate. By the above
argumentation, computers cannot (yet) understand this language, or any language
which deviates from the input on which they were trained. I anticipate that the
proportion of human language which can be understood and produced by computers
will grow continually, eventually approaching full human competence, but I hold
that to fully accomplish this task, they will need bodies capable of performing
all the same sensorimotor activities as humans.\footnote{Eventually, they would
become full philosophical zombies, as discussed by \citet{Chalmers96}.}

Thus, I argue that computers are only capable of producing and understanding
human language insofar as they are capable of becoming indistinguishable from
humans, but that this is fully possible in principle. I despair at computers
learning language as they are currently taught, however. It is only by being
embodied and living socially that any being, synthetic or otherwise, can hope
to be a fully literate being. I, for one, welcome our new android
\st{overlords} companions.

\clearpage
\bibliography{course,extra}

\end{document}
